% =====================================================================
% chapters/04_Bench.tex — rendu LaTeX de ../traductions/04_bench.en_brut.txt (23 septembre 2026)
% Section 4 de la VERSION FRANÇAISE de l'article court. \input depuis main.tex, après 03_Agent.tex.
% Ce fichier DÉFINIT \label{sec:bench}, visé par la carte de lecture de la section 1.4.
% Tableaux : tab:scale en table* sur les deux colonnes, légende AU-DESSUS, booktabs.
% Équation hors texte : une, non numérotée (\[ ... \]).
% Citations : horl2021eqasim, tisseo2023emc2.
% Note de bas de page : PROGEDO-Diffusion.
% Rendu mot pour mot de la traduction fournie par l'auteur.
% =====================================================================

\section{Banc d'essai}
\label{sec:bench}

Nous comparons quinze décideurs au sein d'une même cohorte, selon un protocole unique, avec un ensemble de mesures.

\subsection{La cohorte}
\label{sec:cohort}

Chaque score présenté dans cet article provient d'une cohorte unique et sécurisée de 1\,000 personas synthétiques. Un persona est une personne synthétique possédant un foyer et un programme de déplacements journaliers. Ces personas forment 499 foyers complets, répartis dans les 453 communes de notre référence. La chaîne de synthèse eqasim \citep{horl2021eqasim} a permis de constituer le panel de 11\,329 personnes dont ils sont issus. Cette référence est l'enquête sur les déplacements des ménages de Toulouse \citep{tisseo2023emc2}, dans laquelle les résidents décrivent chaque déplacement de la veille.

Treize marges de la cohorte sont comparées aux données de l'enquête. Toutes les treize sont conformes. Chaque marge représente la part d'un trait, tel que la tranche d'âge ou la possession d'une voiture. Le seuil d'équivalence était d'un point, fixé avant la mesure. L'écart maximal est de 0,50 point (Annexe A).

La cohorte effectue un peu moins de déplacements que la population cible de l'enquête : 3,30 déplacements par personne et 3,69 par personne mobile, contre 3,53 et 3,95 respectivement. La journée évaluée compte 3\,299 déplacements.

Nous publions le protocole, le code, les données initiales et la cohorte scellée. Notre convention de recherche interdit la diffusion des microdonnées de l'enquête. L'enquête est disponible sur demande individuelle auprès des archives PROGEDO-Diffusion\footnote{\url{https://www.progedo.fr/donnees-disponibles/donnees-localisees/}}.

\subsection{Le protocole en trois règles}
\label{sec:rules}

Trois règles permettent d'éviter trois biais dans la comparaison entre un agent génératif et un modèle tabulaire.

La règle 1 attribue à chaque décideur les mêmes 21 variables : douze pour la personne et le ménage, trois pour le déplacement et six pour la géométrie (Annexe B).

La règle 2 attribue une probabilité à chaque option, et non à l'option qu'il classe en premier. Un classement seul placerait toutes les personnes d'un même profil dans le même mode.

La règle 3 restreint chaque prédiction tabulaire aux modes de transport réellement proposés pour ce trajet, puis la ramène à 100\,\%.

Nous égalisons les 21 variables, et non les informations détenues par chaque partie. L'agent reçoit trois éléments qu'aucun modèle tabulaire ne peut recevoir : le programme de ses trajets restants, la météo des prochaines heures et chaque étape de l'itinéraire. Les modèles tabulaires de référence ont reçu un élément que l'agent n'a jamais reçu : les 39\,203 trajets de l'enquête sur lesquels ils ont été estimés. L'asymétrie d'exposition s'exerce dans l'autre sens, ce qui confère à ces modèles une valeur de référence élevée.

\subsection{Ce que nous mesurons}
\label{sec:quantities}

Nous mesurons à deux échelles : la répartition modale de la cohorte et la concordance pour chaque trajet. La répartition modale correspond à la part des trajets effectués en voiture, en transports en commun, à pied et à vélo. L'enquête fournit la part de référence au niveau global, puis au sein de cinq strates. Trois sont nominales : profession, sexe et motif du trajet. Deux sont ordinales : classe d'âge et classe de distance.

Un nombre composite représente la valeur agrégée. Il additionne la divergence globale à la divergence moyenne pondérée de chaque strate.
\[
  \mathcal{C}_{\text{EMD--JSD}} = \mathrm{JSD}^{\text{global}}
  + \sum_{d\ \text{nominal}} w_d\,\overline{\mathrm{JSD}}^{\,d}
  + \sum_{d\ \text{ordinal}} w_d\,\overline{\mathrm{EMD}}^{\,d}
\]
La divergence est celle de Jensen-Shannon pour les strates nominales et la distance de Wasserstein pour les strates ordonnées, respectant ainsi l'ordre des classes. Un composite plus faible indique une distribution plus proche de la distribution de référence.

Deux autres valeurs accompagnent le composite ; les trois sont publiées ensemble. La première recalcule le composite sans les trajets n'offrant qu'une seule option, pour lesquels aucune préférence n'est exprimée. La seconde est l'erreur L1 sur les parts globales, soit la somme des écarts absolus en points de pourcentage.

Une différence entre deux décideurs n'est prise en compte que si elle dépasse l'impact qu'aurait une nouvelle cohorte. Une autre cohorte de 1\,000 profils, construite de la même manière, modifierait un composite de $\pm$1,3 point. La résolution varie selon le décideur, de 0,9 point pour les conditions optimisées et tabulaires à 2,1 pour les conditions minimales. Les déplacements d'une même personne n'étant pas indépendants, chaque intervalle résulte d'un rééchantillonnage des groupes au niveau individuel. Une différence appariée compare deux décideurs sur les personnes qu'ils ont tous deux évaluées, à 95\,\%.

L'audit au niveau unitaire pose l'autre question, sur les jours réellement décrits par les répondants. Il couvre 9\,621 déplacements de 2\,930 répondants, chacun correspondant à son mode déclaré. Nous publions l'exactitude, l'entropie croisée, le rappel et la précision par mode. L'entropie croisée diminue lorsqu'un décideur accorde plus de probabilité au mode déclaré. Le rappel correspond à la part des trajets déclarés d'un mode qu'il désigne.

\subsection{Niveaux, références, décideurs}
\label{sec:deciders}

Trois niveaux définissent le niveau atteint par un décideur sans connaissance comportementale. Le premier niveau répartit uniformément les options proposées, le deuxième privilégie la voiture, le mode majoritaire de la zone. Le troisième, l'heuristique de durée minimale, conserve l'itinéraire le plus rapide proposé.

Quatre références tabulaires indiquent le niveau atteint par l'estimation sur l'enquête. Il s'agit d'un modèle logit multinomial, du modèle de choix discret standard, ainsi que d'un modèle de gradient boosting, d'une régression logistique à noyau et d'une forêt aléatoire. Nous estimons ces quatre modèles à parité stricte, sur la même partition de l'enquête. Aucun ne surpasse les trois autres. Nous lisons donc le plafond une lecture à la fois.

Trois familles de décideurs sont testées. L'invite minimale fournit à un modèle de langage les faits du voyage et le format de sortie, sans critère général. L'invite experte conserve ce format et attire l'attention du modèle sur quatre critères généraux. Deux concernent le confort : la marche et l'attente des transports en commun sont importantes, et un rythme tranquille est recommandé pour un voyageur âgé. Les deux autres concernent les contraintes : des sacs de courses lourds à porter et une journée de travail sans temps mort. La troisième famille est un classificateur zéro-shot avec sortie typée (TypeSafe System One, \texttt{jev-1.13.0}). Il lit ce même texte et renvoie une probabilité par option sans écrire de phrase.

Nous avons obtenu l'invite experte par mutations successives d'un texte, sous trois contraintes. Chaque itération lit les lacunes par strate de l'invite en service, propose une réécriture ciblée et la conserve lorsque le mode surreprésenté disparaît. Les contraintes interdisent un seuil numérique, la dénomination d'une catégorie d'enquête et exigent une notation basée sur des distributions verbalisées.

Le classificateur dispose de sa propre invite d'expert, optimisée pour la première cohorte. Cette invite ajoute un cinquième critère général : l'exposition réelle aux conditions météorologiques. Chaque croisement de modèles et d'invites est donc mesuré sur une seconde cohorte, qui ne partage aucune caractéristique avec la première.

\begin{table*}[tp]
  \caption{Les quinze décideurs lors des trois lectures, avec l'intervalle inter-semences où il a été mesuré. Plus le score est bas, meilleur il est sur les trois lectures.}
  \label{tab:scale}
  \begin{tabular}{lrrrl}
    \toprule
    Décideur & Composite & Excluant l'option unique & L1 sur les parts globales & Intervalle inter-semences \\
    \midrule
    Aléatoire uniforme                                     & 50,16 & 57,90 & 86,80 & déterministe \\
    Tout-voiture                                           & 30,73 & 28,41 & 58,00 & déterministe \\
    Durée minimale                                         & 26,97 & 23,93 & 53,62 & déterministe \\
    Invite minimale, \texttt{mistral-large}                & 14,75 & 20,72 & 44,71 & non rejoué \\
    Invite minimale, classificateur typé                   & 13,75 & 19,84 & 42,76 & non rejoué \\
    Invite minimale, \texttt{gemini-3.1}                   & 12,38 & 16,65 & 39,88 & non rejoué \\
    Invite minimale, \texttt{gemini-3.5}                   &  7,02 & 10,39 & 24,08 & \textbf{[en attente]} \\
    Invite experte, \texttt{gemini-3.1}                    &  8,98 & 12,39 & 31,25 & non rejoué \\
    Invite experte, \texttt{mistral-large}                 &  7,63 & 12,27 & 26,64 & non rejoué \\
    Invite experte, \texttt{gemini-3.5}                    &  4,86 &  6,86 & 13,85 & 0,56 \\
    Logit multinomial                                      &  4,02 &  6,63 &  8,99 & déterministe \\
    Forêt aléatoire                                        &  4,09 &  5,81 &  5,28 & déterministe \\
    Invite experte, classificateur typé (dans l'échantillon) &  3,65 &  6,58 & 10,48 & 0,05 \\
    Régression logistique à noyau                          &  3,61 &  5,61 &  6,69 & déterministe \\
    Gradient boosting                                      &  3,60 &  5,84 &  9,49 & déterministe \\
    \bottomrule
  \end{tabular}
\end{table*}

Chaque score suivant est issu de ce banc d'essai, en commençant par la journée ordinaire décrite par l'enquête.
